\section{Consequences}
\label{sec:applications}

No new quantum calculation is needed in this section.  We first apply the
invariant to products of threefolds with \(\PP^1\), then transfer the stabilization theorem
through Kuznetsov's flop, and finally combine it with known criteria for
universal \(CH_0\)-triviality.

\begin{corollary}[Threefold criterion]
\label{cor:threefold-marker}
\coverage{fragment}%
\lean{threefold_not_rational_of_occurrenceIndexedMarker}%
Let \(Y\) be a smooth projective complex threefold.  If its generic even QDM
contains a block counted by the exponent count defined in
\eqref{eq:atomic-fold}, then \(Y\times\PP^1\) is irrational.  In particular,
\(Y\) is irrational.
\end{corollary}

\begin{proof}[Proof of Corollary~\ref{cor:threefold-marker}]
A counted block gives \(I_{\mathrm{exp}}(Y)>0\).  The projective-bundle
formula gives
\[
 I_{\mathrm{exp}}(Y\times\PP^1)=2I_{\mathrm{exp}}(Y)>0.
\]
Every center in projective weak factorization of fourfolds has dimension
at most two, so Proposition~\ref{prop:atomic-lowdim} makes every correction
zero.  Theorem~\ref{thm:marker-ledger} therefore makes the invariant
birationally invariant in dimension four.  Since
\(I_{\mathrm{exp}}(\PP^4)=0\) by \eqref{eq:atomic-projective-value},
\(Y\times\PP^1\) is irrational.  Rationality of \(Y\) would imply
rationality of this product, so \(Y\) is irrational as well.
\proves{cor:threefold-marker}%
\uses{prop:atomic-lowdim, thm:marker-ledger, prop:qdm-operation-ledgers}%
\end{proof}

Kuznetsov's associated Pfaffian cubic and flop supply the next application.
Tschinkel and Zhang study equivariant and twisted forms of this correspondence
\cite{TschinkelZhangV14}.

\begin{corollary}[Irrationality after one stabilization for \(V_{14}\)]
\label{cor:v14-one-step}
\coverage{conditional_deduction}%
\lean{genusEight_oneProjectiveLine_not_rational_of_residueContext}%
For every smooth prime Fano threefold \(V\) of genus eight,
\(V\times\PP^1\) is irrational.
\end{corollary}

\begin{proof}[Proof of Corollary~\ref{cor:v14-one-step}]
Kuznetsov's flop gives a birational map between rank-two projective bundles
\(\PP_V(U)\) and \(\PP_X(E^*)\) over \(V\) and its associated Pfaffian cubic
threefold \(X\) \cite[Theorems~2.17--2.18]{KuznetsovV14}.  Each projective
bundle is birational to the product of its base with \(\PP^1\).  Hence
\(V\times\PP^1\) is birational to \(X\times\PP^1\), which is irrational by
Theorem~\ref{thm:every-cubic}.
\proves{cor:v14-one-step}%
\uses{thm:every-cubic}%
\imports{KuznetsovV14:flop}%
\end{proof}

\begin{corollary}[Separation in codimension at most three]
\label{cor:voisin-separation}
\coverage{fragment}%
\lean{primitiveMinimalClassFamily_universalCH0_and_irrational_stabilization,%
 exists_universalCH0_cubic_with_irrational_stabilization}%
\uses{thm:every-cubic}%
\imports{Voisin:minimal-class-locus}%
There is a nonempty countable union of subvarieties of codimension at most
three in the moduli space of smooth cubic threefolds such that, for every
member \(X\), both \(X\) and \(X\times\PP^1\) are universally
\(CH_0\)-trivial, while \(X\times\PP^1\) is irrational.
\end{corollary}

\begin{proof}[Proof of Corollary~\ref{cor:voisin-separation}]
Voisin produces a nonempty countable union of subvarieties of codimension at
most three in the moduli space of smooth cubic threefolds along which the
primitive minimal class \(\Theta^4/4!\) is algebraic, hence along which
\(CH_0\) is universally trivial \cite[Theorem~4.5 and
Lemma~4.6]{Voisin}.  The projective-bundle formula passes universal
\(CH_0\)-triviality to the product, and
Theorem~\ref{thm:every-cubic} supplies its irrationality.
\proves{cor:voisin-separation}%
\uses{thm:every-cubic}%
\imports{Voisin:minimal-class-locus}%
\end{proof}

\begin{corollary}[Fermat cubic]
\label{cor:fermat-separation}
\coverage{conditional_deduction}%
\lean{fermatCubic_universalCH0_and_irrational_stabilization}%
\uses{thm:every-cubic}%
\imports{CT:separated-variable-criterion}%
Let
\[
  X=\{x_1^3+x_2^3+x_3^3+x_4^3+x_5^3=0\}\subset\PP^4.
\]
Then the fourfold \(X\times\PP^1\) is universally \(CH_0\)-trivial and
irrational.
\end{corollary}

\begin{proof}[Proof of Corollary~\ref{cor:fermat-separation}]
The Fermat equation is a sum of five forms in separated single variables, so
\cite[Theorem~2.8]{CT} gives universal \(CH_0\)-triviality; the projective-bundle formula
passes it to the product.  Irrationality follows from
Theorem~\ref{thm:every-cubic}.
\proves{cor:fermat-separation}%
\uses{thm:every-cubic}%
\imports{CT:separated-variable-criterion}%
\end{proof}

\begin{corollary}[Coprime-degree cubics]
\label{cor:coprime-separation}
\coverage{conditional_deduction}%
\lean{coprimeUnirationalDegrees_universalCH0_and_irrational_stabilization}%
\uses{thm:every-cubic}%
\imports{YYZ:coprime-degree-family, YYZ:coprime-degrees-persist}%
Let \(X\) be one of the cubic threefolds of
\cite[Theorem~3.3]{YYZ}.  Then \(X\times\PP^1\) admits unirational
parametrizations of coprime degrees, is universally \(CH_0\)-trivial, and is
irrational.
\end{corollary}

\begin{proof}[Proof of Corollary~\ref{cor:coprime-separation}]
The cubics of \cite[Theorem~3.3]{YYZ} carry unirational parametrizations of
degrees two and three, and \cite[Corollary~3.5]{YYZ} keeps both degrees on
\(X\times\PP^1\).  Parametrizations of coprime degrees give a decomposition
of the diagonal and hence universal \(CH_0\)-triviality
\cite[Section~1]{YYZ}.  Theorem~\ref{thm:every-cubic} gives irrationality.
\proves{cor:coprime-separation}%
\uses{thm:every-cubic}%
\imports{YYZ:coprime-degree-family, YYZ:coprime-degrees-persist}%
\end{proof}

In dimension five, threefold centers can contribute to weak factorization.
The next section retains those contributions and places the blowup formulas
in the additive Grothendieck group.
